\section{Absolute continuity and singularity}
\label{sec:abscont}

Throughout, $(\Xset, \salg)$ is a measurable space and all measures are defined
on $\salg$. A measure $\mu$ is \emph{$\sigma$-finite} if $\Xset = \bigcup_n A_n$
with $A_n \in \salg$ and $\mu(A_n) < \infty$; $\sigma$-finiteness is the standing
hypothesis under which the theorems below hold in full generality.

\begin{definition}[Absolute continuity]
\label{def:abscont}
Let $\mu, \nu$ be measures on $(\Xset, \salg)$. We say $\nu$ is
\emph{absolutely continuous} with respect to $\mu$, written $\nu \abscont \mu$,
if
\begin{equation}
  \mu(A) = 0 \;\Longrightarrow\; \nu(A) = 0
  \qquad \text{for every } A \in \salg.
\end{equation}
Equivalently: every $\mu$-null set is $\nu$-null.
\end{definition}

\begin{definition}[Mutual singularity]
\label{def:singular}
The measures $\mu, \nu$ are \emph{mutually singular}, written $\mu \sing \nu$, if
there is a set $E \in \salg$ with
\begin{equation}
  \mu(E) = 0 \qquad \text{and} \qquad \nu(\Xset \setminus E) = 0,
\end{equation}
i.e.\ they live on disjoint sets: $\nu$ is carried by $E$ and $\mu$ by its
complement.
\end{definition}

Absolute continuity and singularity are the two extremes. The Lebesgue
decomposition of Section~\ref{sec:lebesgue} shows every $\sigma$-finite $\nu$
splits uniquely into a piece of each type relative to $\mu$.

\paragraph{The \texorpdfstring{$\varepsilon$--$\delta$}{epsilon-delta} form.}
For \emph{finite} $\nu$, the set-theoretic Definition~\ref{def:abscont} is
equivalent to an analytic continuity statement, which is where the name comes
from.

\begin{proposition}[$\varepsilon$--$\delta$ characterization]
\label{prop:eps-delta}
Let $\nu$ be a finite measure. Then $\nu \abscont \mu$ if and only if for every
$\varepsilon > 0$ there exists $\delta > 0$ such that
\begin{equation}
  A \in \salg,\ \mu(A) < \delta \;\Longrightarrow\; \nu(A) < \varepsilon.
  \label{eq:eps-delta}
\end{equation}
\end{proposition}

\begin{proof}[Proof sketch]
The $\varepsilon$--$\delta$ condition trivially forces $\mu(A)=0 \Rightarrow
\nu(A) = 0$. Conversely, if \eqref{eq:eps-delta} failed there would be
$\varepsilon_0 > 0$ and sets $A_n$ with $\mu(A_n) < 2^{-n}$ but
$\nu(A_n) \ge \varepsilon_0$. Put $A = \limsup_n A_n$. By Borel--Cantelli on
$\mu$, $\mu(A) = 0$, while $\nu(A) \ge \varepsilon_0$ by finiteness of $\nu$ and
reverse Fatou, contradicting $\nu \abscont \mu$.
\end{proof}

\begin{remark}
Finiteness matters in Proposition~\ref{prop:eps-delta}. With
$\mu = $ Lebesgue measure and $\nu(A) = \int_A x^2\,dx$ on $\R$, one has
$\nu \abscont \mu$ (Definition~\ref{def:abscont}) yet no single $\delta$ works
uniformly for the interval family $\{[n, n+\delta]\}$, because $\nu$ is not
finite. The remedy is to test \eqref{eq:eps-delta} on sets of finite $\nu$-measure.
\end{remark}

\paragraph{Signed measures and the Jordan decomposition.}
A \emph{signed measure} $\nu$ assigns values in $(-\infty, +\infty]$ (or
$[-\infty, +\infty)$) and is countably additive. The Hahn decomposition gives a
partition $\Xset = P \sqcup N$ into a positive set $P$ and a negative set $N$;
setting $\nu^+(A) = \nu(A \cap P)$ and $\nu^-(A) = -\nu(A \cap N)$ yields two
genuine (nonnegative) measures.

\begin{proposition}[Jordan decomposition]
\label{prop:jordan}
Every signed measure $\nu$ admits a unique decomposition
\begin{equation}
  \nu = \nu^+ - \nu^-, \qquad \nu^+ \sing \nu^-,
\end{equation}
into mutually singular measures. Its \emph{total variation} is
$\lvert \nu \rvert = \nu^+ + \nu^-$, and $\nu \abscont \mu$ is defined to mean
$\lvert \nu \rvert \abscont \mu$.
\end{proposition}

This reduces every statement about signed measures to the nonnegative case
applied to $\nu^+$ and $\nu^-$ separately, which is how the Radon--Nikodym
theorem is extended to signed and complex measures.
